Przypomnijmy pojęcia diagramu i stożka wprowadzone pod koniec poprzedniego rozdziału.

\begin{definicja}
Niech $\Jcat$ i $\C$ będą kategoriami. \emph{Diagramem typu $\Jcat$ w $\C$} nazywamy funktor
\[
D \colon \Jcat \longrightarrow \C.
\]
\end{definicja}

\begin{uwaga}
Kategorię $\Jcat$ nazywamy \emph{kategorią indeksów}; jej obiekty oznaczamy zwykle przez $i, j, k, \dots$, a morfizmy przez $\alpha, \beta, \gamma, \dots$. Wartości funktora $D$ oznaczamy odpowiednio przez $D_i, D_j, \dots$ oraz $D_\alpha, D_\beta, \dots$
\end{uwaga}

\begin{definicja}
\emph{Stożkiem nad diagramem $D \colon \Jcat \to \C$} nazywamy obiekt $C \in \C$ wraz z rodziną morfizmów
\[
\bigl\{\, c_i \colon C \longrightarrow D_i \,\bigr\}_{i \in \Ob \Jcat},
\]
takich, że dla każdego $\alpha \colon i \to j$ w $\Jcat$ diagram
\[
\begin{tikzcd}[row sep=large, column sep=large]
 & C \arrow[dl, "c_i"'] \arrow[dr, "c_j"] & \\
 D_i \arrow[rr, "D_\alpha"'] & & D_j
\end{tikzcd}
\]
jest przemienny.
\end{definicja}

\begin{definicja}
Niech $(C, \{c_i\})$ oraz $(C', \{c'_i\})$ będą stożkami nad diagramem $D \colon \Jcat \to \C$. \emph{Morfizmem stożków} $(C, \{c_i\}) \to (C', \{c'_i\})$ nazywamy strzałkę
\[
\vartheta \colon C \longrightarrow C' \quad \text{w } \C,
\]
taką, że diagram
\[
\begin{tikzcd}[row sep=large, column sep=large]
C \arrow[rr, "\vartheta"] \arrow[dr, "c_i"'] & & C' \arrow[dl, "c'_i"] \\
 & D_i &
\end{tikzcd}
\]
jest przemienny dla każdego $i \in \Ob \Jcat$.
\end{definicja}

\begin{definicja}
Przez $\mathrm{Cone}(D)$ oznaczamy kategorię, w której obiektami są stożki nad diagramem $D \colon \Jcat \to \C$, a morfizmami --- morfizmy stożków.
\end{definicja}

\begin{definicja}
\emph{Granicą diagramu $D \colon \Jcat \to \C$} nazywamy obiekt końcowy w kategorii $\mathrm{Cone}(D)$ i oznaczamy go przez
\[
\Bigl( \varprojlim_{j} D_j,\ \bigl\{ p_i \colon \varprojlim_{j} D_j \to D_i \bigr\} \Bigr).
\]
\end{definicja}

\subsection*{Przykłady}

\begin{przyklad}
Jeśli $\Jcat = \emptyset$, to $\mathrm{Cone}(D) = \C$, więc
\[
\varprojlim D \;=\; \text{obiekt końcowy w } \C.
\]
\end{przyklad}

\begin{przyklad}
Jeśli $\Jcat$ jest kategorią dyskretną o dwóch obiektach $\bullet_1, \bullet_2$, a $D(\bullet_i) = D_i$, to
\[
\varprojlim D \;=\; D_1 \times D_2.
\]
\end{przyklad}

\begin{przyklad}
Jeśli $\Jcat$ jest kategorią z dwoma obiektami i dwiema równoległymi strzałkami między nimi:
\[
\Jcat \;=\; \begin{tikzcd} \bullet_1 \arrow[r, shift left, "\alpha"] \arrow[r, shift right, "\beta"'] & \bullet_2 \end{tikzcd},
\qquad
D \;=\; \begin{tikzcd} D_1 \arrow[r, shift left, "D_\alpha"] \arrow[r, shift right, "D_\beta"'] & D_2 \end{tikzcd},
\]
to
\[
\varprojlim D \;=\; \text{ekwalizator } D_\alpha,\, D_\beta.
\]
\end{przyklad}

\begin{przyklad}
Jeśli $\Jcat$ jest kategorią postaci
\[
\Jcat \;=\; \begin{tikzcd} & \bullet_3 \arrow[d, "\beta"] \\ \bullet_1 \arrow[r, "\alpha"'] & \bullet_2 \end{tikzcd},
\qquad
D \;=\; \begin{tikzcd} & D_3 \arrow[d, "D_\beta"] \\ D_1 \arrow[r, "D_\alpha"'] & D_2 \end{tikzcd},
\]
to
\[
\varprojlim D \;=\; \text{pullback } D_\alpha, D_\beta \text{ w } \C.
\]
\end{przyklad}

\begin{przyklad}[Granica nieskończonego łańcucha]
Granice istnieją także nad diagramami nieskończonymi. Niech $\Jcat$ będzie kategorią indukowaną przez poset $(\mathbb{N}, \ge)$; diagram $D\colon \Jcat\to\Set$ to łańcuch zbiorów i przekształceń
\[
\cdots \longrightarrow X_2 \xrightarrow{\;f_1\;} X_1 \xrightarrow{\;f_0\;} X_0.
\]
W $\Set$ granicą takiego diagramu (zwaną \emph{granicą odwrotną}) jest
\[
\varprojlim_n X_n \;=\; \Bigl\{ (x_0, x_1, x_2, \ldots) \in \prod_n X_n \;:\; f_n(x_{n+1}) = x_n \text{ dla każdego } n \Bigr\}
\]
wraz z rzutowaniami. Na przykład dla $X_n = A^n$ (ciągi długości $n$ nad zbiorem $A$) i przekształceń obcinających ostatni element granicą jest zbiór $A^\omega$ wszystkich nieskończonych ciągów (strumieni) nad $A$.
\end{przyklad}

\begin{definicja}
\emph{Skończoną granicą} nazywamy granicę nad diagramem $D \colon \Jcat \to \C$, w którym kategoria $\Jcat$ jest skończona.
\end{definicja}

\begin{twierdzenie}
Kategoria $\C$ ma wszystkie skończone granice wtedy i tylko wtedy, gdy ma skończone produkty i ekwalizatory (lub równoważnie: gdy ma pullbacki i obiekt końcowy).
\end{twierdzenie}

\begin{proof}[Dowód (szkic)]
Pokażemy, że z faktu istnienia skończonych produktów i ekwalizatorów wynika istnienie wszystkich skończonych granic.

Weźmy diagram $D \colon \Jcat \to \C$, gdzie $\Jcat$ jest skończoną kategorią indeksów. Podamy konstrukcję obiektu $E$ wraz ze strzałkami $e_i \colon E \to D_i$, dla którego $(E, \{e_i\})$ będzie obiektem końcowym w $\mathrm{Cone}(D)$.

\textbf{(1)} Rozważmy następujące obiekty w $\C$:
\[
\prod_{i \in \Ob \Jcat} D_i \qquad \text{oraz} \qquad \prod_{\alpha \colon i \to j \in \Mor \Jcat} D_j.
\]

\textbf{(2)} Zdefiniujmy dwie strzałki
\[
\begin{tikzcd}[column sep=large]
\displaystyle\prod_{i \in \Ob \Jcat} D_i \arrow[r, shift left, "\varphi"] \arrow[r, shift right, "\psi"'] & \displaystyle\prod_{\alpha \colon i \to j} D_j,
\end{tikzcd}
\]
zadane następująco: rzutowanie $\pi_\alpha$ z produktu po prawej stronie spełnia
\[
\pi_\alpha \circ \varphi = \pi_j, \qquad \pi_\alpha \circ \psi = D_\alpha \circ \pi_i.
\]

\textbf{(3)} Weźmy ekwalizator $E$ wraz ze strzałką $e \colon E \to \prod_{i \in \Ob \Jcat} D_i$ pary $\varphi, \psi$:
\[
\begin{tikzcd}[column sep=large]
E \arrow[r, "e"] & \displaystyle\prod_{i \in \Ob \Jcat} D_i \arrow[r, shift left, "\varphi"] \arrow[r, shift right, "\psi"'] & \displaystyle\prod_{\alpha \colon i \to j} D_j.
\end{tikzcd}
\]
Definiujemy $e_i := \pi_i \circ e$. Można pokazać, że $E$ wraz z $\{e_i\}_{i \in \Ob \Jcat}$ spełnia żądane własności.
\end{proof}

\subsection*{Zachowywanie granic}

\begin{definicja}
Mówimy, że funktor $F \colon \C \to \D$ \emph{zachowuje granice} nad diagramem typu $\Jcat$, jeśli dla każdego diagramu $D \colon \Jcat \to \C$ zachodzi
\[
F\Bigl(\varprojlim_i D_i\Bigr) \;\cong\; \varprojlim_i F(D_i)
\]
(izomorfizm kanoniczny).

Mówimy, że funktor $F \colon \C \to \D$ \emph{zachowuje wszystkie granice}, jeśli zachowuje granice nad diagramami każdego typu.
\end{definicja}

\begin{przyklad}
Niech $\C$ będzie kategorią. Rozważmy funktor
\[
\C(X, -) \colon \C \longrightarrow \Set
\]
dla ustalonego obiektu $X \in \C$:
\[
A \longmapsto \C(X, A) = \text{zbiór strzałek } X \to A \text{ w } \C,
\]
oraz dla $f \colon A \to B$:
\[
\C(X, A) \xrightarrow{\;\C(X, f)\;} \C(X, B), \qquad \C(X, f)(g \colon X \to A) \;=\; f \circ g.
\]
(Często $\C(X, f) =: f_*$, czyli $f_*(g) = f \circ g$.)

Pokażemy, że funktor $\C(X, -)$ zachowuje wszystkie skończone granice.

\textbf{(1) Obiekt końcowy.} Jeśli $\mathbf{1}$ jest obiektem końcowym w $\C$, to
\[
\C(X, \mathbf{1}) \;=\; \{!_X\} \quad \text{--- jednoelementowy zbiór,}
\]
czyli obiekt końcowy w $\Set$.

\textbf{(2) Produkty.} Z ćwiczeń (zadanie~\ref{zad:hom-produkty}) wiemy, że $\C(X, -)$ zachowuje produkty.

\textbf{(3) Ekwalizatory.} Weźmy parę $f, g \colon A \rightrightarrows B$, ich ekwalizator
\[
\begin{tikzcd}
E \arrow[r, "e"] & A \arrow[r, shift left, "f"] \arrow[r, shift right, "g"'] & B.
\end{tikzcd}
\]
Czy diagram
\[
\begin{tikzcd}
\C(X, E) \arrow[r, "e_*"] & \C(X, A) \arrow[r, shift left, "f_*"] \arrow[r, shift right, "g_*"'] & \C(X, B)
\end{tikzcd}
\]
jest diagramem ekwalizatora w $\Set$? Weźmy $x \in \C(X, A)$, czyli $x \colon X \to A$. Warunek $f_*(x) = g_*(x)$ jest równoważny $f \circ x = g \circ x$. Z własności ekwalizatora istnieje dokładnie jedna strzałka $u_x \colon X \to E$ taka, że $e \circ u_x = x$. Czyli
\[
e_* \colon \C(X, E) \longrightarrow \{\, x \in \C(X, A) \colon f_*(x) = g_*(x) \,\}
\]
jest bijekcją, a więc $\C(X, E)$ jest ekwalizatorem $f_*, g_*$.
\end{przyklad}

\begin{wniosek}
Funktor $\C(X, -)$ zachowuje wszystkie skończone granice.
\end{wniosek}

\begin{uwaga}
Można pokazać, że $\C(X, -)$ zachowuje wszystkie granice (nie tylko skończone).
\end{uwaga}

\begin{przyklad}[Funktor, który nie zachowuje granic]
Nie każdy funktor zachowuje granice. Rozważmy $F = (-) + 1 \colon \Set \to \Set$, $FX = X + \{\ast\}$. Funktor $F$ nie zachowuje obiektu końcowego: $F(\mathbf{1})$ jest zbiorem dwuelementowym, a więc nie jest obiektem końcowym w $\Set$. Nie zachowuje też produktów: zbiór $F(\mathbf{1}\times\mathbf{1})$ ma dwa elementy, podczas gdy $F(\mathbf{1})\times F(\mathbf{1})$ -- cztery.
\end{przyklad}
